\section{Upper Breadboard Setup and Conveyor Belt Alignment}
\label{sec:upper_breadboard}

The implementation of the optical conveyor belt requires the spatial superposition of two counter-propagating optical fields: the bottom-up beam delivered by the Hollow-Core Photonic Crystal Fiber (HCPCF) through the vacuum chamber, and the top-down beam delivered by the black fiber. The combination and alignment of these two optical paths take place on the upper breadboard, located above the main experimental chamber.

The experimental work on the upper breadboard was divided into three main phases: the outcoupling and conditioning of the beam emerging from the black fiber, the design and construction of dedicated imaging systems to monitor the spatial modes, and the systematic alignment of the counter-propagating beams to maximize their spatial overlap.

\subsection{Black-fiber outcoupling and dual-purpose injection path}

The light injected into the black fiber on the lower bench (as described in Section~\ref{sec:injection_black_fiber}) is guided to the upper breadboard. To ensure pointing stability and minimize mechanical drifts, the outcoupling facet of the black fiber was secured using a custom-made support, which was shortened by $3\,\mathrm{mm}$ compared to its original design to increase mechanical rigidity.

The diverging beam emerging from the fiber was collimated using an $f=19\,\mathrm{mm}$ lens, positioned manually using a high-stability mount. Following the collimation stage, a beam expander providing a magnification of approximately $3\times$ was installed and aligned. The purpose of this telescope is to match the transverse size of the black-fiber beam to the expected size of the beam emerging from the HCPCF, which is a fundamental prerequisite for achieving a high-contrast interference pattern in the conveyor-belt configuration. 

The expanded beam is then steered by two mirrors towards a dichroic mirror, which reflects the $1064\,\mathrm{nm}$ light downwards. After the dichroic mirror, the beam passes through an $f=200\,\mathrm{mm}$ lens mounted along the vertical axis. This lens fulfills a crucial dual purpose in the setup: it focuses the top-down beam from the black fiber onto the upper tip of the HCPCF, while simultaneously collimating the counter-propagating bottom-up beam emerging from the same hollow-core fiber. 

Because the optical components on the upper breadboard affect the polarization state of the light, additional waveplates and polarizers are required in the definitive layout to ensure that both counter-propagating beams share the same polarization axis.

\begin{figure}[h!]
    \centering
    % Inserisci il nome del file pdf senza spazi
    \includegraphics[width=0.8\textwidth]{Cap_2_figures/upper.pdf} 
    \caption{Questa è un'immagine vettoriale caricata da un PDF. MI sa che l'AOM va ruotato}
    \label{fig:upper_setup}
\end{figure}


\subsection{Beam propagation and constraints of the optical elements}

The proper design of the shared optical path requires a careful evaluation of the beam size. The expected waist radius $w_{\mathrm{col}}$ of the bottom-up collimated beam can be estimated from the nominal mode radius within the hollow core, $w_{\mathrm{HCPCF}} \approx 19\,\mu\mathrm{m}$, using the relation:

\begin{equation}
    w_{\mathrm{col}} = \frac{\lambda f}{\pi w_{\mathrm{HCPCF}}}.
\end{equation}

For $\lambda = 1064\,\mathrm{nm}$ and the $f = 200\,\mathrm{mm}$ vertical lens, the calculated beam radius is $w_{\mathrm{col}} \approx 3.6\,\mathrm{mm}$, corresponding to a full beam diameter $d_{\mathrm{col}} = 2 w_{\mathrm{col}} \approx 7.2\,\mathrm{mm}$. 

This collimated beam is subsequently steered by standard $1$-inch ($25.4\,\mathrm{mm}$) mirrors. At an incidence angle of $45^\circ$, the beam footprint on the mirror surface becomes elliptical, with the major axis $d^*$ given by:

\begin{equation}
    d^* = \frac{d_{\mathrm{col}}}{\cos(45^\circ)} \approx 10.2\,\mathrm{mm}.
\end{equation}

Although this footprint is mathematically smaller than the clear aperture of a $1$-inch mirror, the available clearance is relatively modest. Furthermore, if the effective mode radius in the HCPCF is smaller than the assumed $19\,\mu\mathrm{m}$ (due to the $48\,\mu\mathrm{m}$ core diameter constraints), the collimated beam divergence and the resulting mirror footprint would be even larger. Consequently, careful alignment—verified with an infrared viewer—is critical to prevent clipping and vignetting at the edges of the steering mirrors.

\subsection{Mode imaging and alignment strategy}

To properly align this counter-propagating configuration and ensure optimal mode matching, it was necessary to simultaneously monitor the spatial mode of the HCPCF and optimize the top-down injection. For this purpose, an imaging system (microscope) was assembled on the upper breadboard, exploiting the shared optical path.

A magnetic flippable mirror was inserted between the dichroic mirror and the first steering mirror of the black-fiber branch to pick off the bottom-up beam. The imaging path consists of an $f=300\,\mathrm{mm}$ lens, a Zeiss microscope objective ($16.5\,\mathrm{mm}$, Plan Neofluar $10\times/0.30$), and a final $f=500\,\mathrm{mm}$ focusing lens that projects the image onto a CMOS camera. This multi-stage magnification setup provides high resolution for resolving the spatial distribution of the light exiting the microstructured fiber.


\begin{figure}[h!]
    \centering
    % Inserisci il nome del file pdf senza spazi
    \includegraphics[width=0.8\textwidth]{Cap_2_figures/upper_microscope.pdf} 
    \caption{Questa è un'immagine vettoriale caricata da un PDF. MI sa che l'AOM va ruotato}
    \label{fig:upper_microscope}
\end{figure}


% INSERT FIGURE:
% - Image of the spatial mode from the HCPCF (bottom-up) taken with the upper camera.
% - Specify in the caption if the "Pastel" pseudocolor map was used.

The primary assumption of the alignment method is that the mode in the HCPCF must be properly injected and symmetrically propagated. During the initial observations with the newly constructed microscope, it was noted that the fundamental mode was not propagating symmetrically within the core. By applying gentle mechanical perturbations to the fiber, higher-order propagation modes were intentionally excited. This behavior was systematically exploited to map the spatial boundaries of the hollow core. Once the core region was identified, the fiber position was finely adjusted to ensure that the mode extended symmetrically across the entire available area. This optimization yielded a symmetric fundamental Gaussian mode with an overall bottom-up injection efficiency of approximately $20\% - 25\%$.

With the HCPCF mode optimized, the spatial overlap between the top-down and bottom-up beams was verified. A reference mirror was placed after the dichroic mirror to project the beam profiles onto a distant viewing card, checking the collimation of the HCPCF beam. The procedure was repeated for the black-fiber beam, sampling it after the beam expander and at a far-field plane. Finally, both beams were turned on simultaneously to ensure that they were correctly centered on all lenses and mirrors and perfectly overlapped along the entire shared path. By substituting the magnetic mirror with a beam splitter, it was possible to operate the microscope continuously, exploiting the small transmission of the dichroic mirror at $1064\,\mathrm{nm}$ to simultaneously optimize the injection and monitor the mode in real-time.


\subsection{Figures of merit for the injection efficiency}

To quantitatively evaluate the success of the optical setup and the alignment procedure, four distinct figures of merit were established and continuously monitored:

\begin{enumerate}
    \item \textbf{Bottom-up HCPCF mode quality:} The symmetry and spatial distribution of the beam emerging from the vacuum chamber, evaluated through the upper breadboard microscope.
    \item \textbf{Bottom-up HCPCF injection efficiency:} The ratio of the transmitted power measured after the vacuum chamber to the incident power measured before the lower injection lens.
    \item \textbf{Retro-injection into the black fiber:} The coupling efficiency of the bottom-up HCPCF beam into the black fiber. Since the two beams must be perfectly collinear to create the conveyor belt, optimizing this reverse-coupling signal inherently guarantees the spatial overlap of the two modes on the upper breadboard.
    \item \textbf{Top-down injection into the HCPCF:} The coupling of the beam delivered by the black fiber downwards into the hollow-core fiber. 
\end{enumerate}

To accurately evaluate the fourth figure of merit, a secondary microscope was built on the lower optical bench. By exploiting the small but non-zero transmission of the lower $1064\,\mathrm{nm}$ dichroic mirror, the light exiting the bottom tip of the HCPCF was picked off using a flippable mirror and an additional steering mirror, and then focused onto a camera using an $f=1000\,\mathrm{mm}$ lens. This diagnostic tool allowed for the simultaneous optimization of the top-down injection and the real-time monitoring of the spatial mode traversing the vacuum chamber from top to bottom.

\subsection{RF electronics characterization for intensity control}

% ADD INTRODUCTORY SENTENCE IF NEEDED:
% e.g., The depth of the optical dipole trap and the dynamics of the conveyor belt require precise control over the optical power of both beams.

The optical power in both the main ODT branch and the black-fiber branch is actively controlled using Acousto-Optic Modulators (AOMs). The Radio-Frequency (RF) signals driving these AOMs are generated by a Direct Digital Synthesizer (DDS) board, specifically designated as DDS3.

Prior to integrating the electronics with the optical setup, the DDS outputs were characterized using a spectrum analyzer. The upper channel of the DDS is dedicated to the primary ODT beam, while the lower channel drives the black-fiber branch (ODT black). For both channels, the emitted RF power was measured as a function of the programmed DDS amplitude parameter. 

% INSERT FIGURE:
% - Plot of the RF power (dBm or Watts) vs DDS amplitude for both channels.
% - Add fitting curves if applicable.